\section{The two computer-assisted steps}\label{sec:certificates}

Two inequalities between explicit real numbers are established by computation: the bound on
the functional $\Psi_0$ behind the integer constant (Theorems~\ref{thm:iid}
and~\ref{thm:integer}) and the bound on the cost of a fixed point of $\RDE$
(Theorem~\ref{thm:existence}). This section states, for each, the inequality that is
certified, the scheme, which parts are computed in interval arithmetic and which in floating
point, how to reproduce the run, and which independent recomputations were made. File paths
are relative to the project folder. The scripts are in Python~3 and use \texttt{mpmath}
version 1.3.0 (pure-Python backend) and, for the second certificate, \texttt{numpy}.

\subsection{The trusted base}\label{sec:certificates-base}

Both certificates rest on the following and on nothing else in the numerical libraries.
\begin{enumerate}
\item[(B1)] The interval type \texttt{mpmath.iv}: addition, subtraction, multiplication,
division and comparison of intervals, and the conversion of integers and binary
floating-point numbers to intervals. In \texttt{mpmath} these are computed from exact integer
arithmetic with directed rounding.
\item[(B2)] Python integers and the exact rationals of \texttt{fractions.Fraction}.
\item[(B3)] For the second certificate only: IEEE~754 binary64 arithmetic, rounding to
nearest, for the additions and multiplications inside sums of products of nonnegative
numbers (computed by \texttt{numpy} and the linear-algebra library behind it).
\item[(B4)] The logic of the scripts.
\end{enumerate}
The following library functions are \emph{not} used: the interval exponential
\texttt{iv.exp}, the interval logarithm \texttt{iv.log}, the interval Euler constant
\texttt{iv.euler}, interval powers, and the printing routine \texttt{iv.nstr}. The reason is
given next. The four interval operations of (B1) were tested against exact rational
arithmetic ($468{,}775$ containment checks at $80$ bits and $160{,}000$ at $120$ bits, no
failure). This is evidence for (B1), not a proof of it.

\subsection{The library's interval exponential and logarithm are not strict enclosures}
\label{sec:certificates-defect}

In \texttt{mpmath} 1.3.0 the interval exponential and logarithm can return an interval that
does not contain the true value.
\begin{itemize}
\item At $80$ bits the upper endpoint of \texttt{iv.exp} was below the true value for $21$ of
$100{,}000$ random arguments, with relative shortfall up to $2.0\cdot10^{-28}$, and the
upper endpoint of \texttt{iv.log} was below the true value for $1$ of $1{,}000{,}000$
arguments, with relative shortfall $3.0\cdot10^{-30}$. The lower endpoints did not fail in
these samples. The reference values were computed at $600$ bits.
\item At $120$ bits the upper endpoint of \texttt{iv.exp} was below a $400$-bit reference
value for $5$ of $20{,}000$ random arguments, by a relative $10^{-40}$.
\item One failure is established without any reference to the library's own elementary
functions. For $x=-\mathtt{0x1.31c2c4b9c3ff0p{-}2}$ the verified exponential described below,
run at $200$ bits, has lower bound $0.741860139211600604812658140023$, while \texttt{iv.exp}
at $80$ bits returns the upper endpoint $0.741860139211600604812658139927$. A failure of
\texttt{iv.log} at $x=\mathtt{0x1.5a6695e45d861p{+}0}$ is established in the same way.
\end{itemize}
The cause, read in the library source (\texttt{libmp/libmpi.py}), is that the interval
routines call the scalar exponential and logarithm with directed rounding, and those compute
an approximation with a few guard bits and round \emph{that} in the requested direction. When
the approximation lies on the wrong side of a representable number, so does the result. The
routine for Euler's constant has the same design. In addition, \texttt{iv.nstr} prints
endpoints rounded to nearest, so a printed ``upper bound'' can be below the endpoint; on
$20{,}000$ random endpoints this happened $10{,}055$ times.

The defect is below $10^{-27}$ in relative size per call and is far too small to change any
digit reported here: the first versions of both certificates used the library functions,
and replacing them left the logs of the second certificate unchanged and moved the bound of
the first by less than $2\cdot10^{-12}$. The certificates were nevertheless changed so that
they do not rest on that estimate. Both now use the following replacements, built from (B1) only.

\emph{Exponential.} Write $x=2^kr$ with $|r|\le2^{-10}$; the inequality is asserted on the
interval endpoints. Then $e^r\in\sum_{j<12}r^j/j!+[-\rho,\rho]$, where $\rho$ is a fixed bound
for the Lagrange remainder $|r|^{12}e^{|r|}/12!$, namely $\rho=2^{-119}/12!$. The sum is
evaluated by Horner's rule in
interval arithmetic, and $e^x$ is obtained by $k$ interval squarings.

\emph{Logarithm} (first certificate only). For a point $p>0$ let $y_0$ be any real number; the
script takes the floating-point logarithm of $p$, which is a guess and is not trusted. Then
$\ln p=y_0+\ln(1+d)$ with $d=p\,e^{-y_0}-1$ enclosed by the verified exponential, and
$d-d^2\le\ln(1+d)\le d$ for $|d|\le1/2$; the script asserts $|d|\le2^{-20}$. An interval is
treated endpoint by endpoint, since the logarithm is increasing. The inequality holds because
$\ln(1+d)\le d$ always, $\ln(1+d)\ge d-d^2/2$ for $d\ge0$, and for $s=-d\in(0,1/2]$,
$-\ln(1-s)=s+s^2/2+s^3/3+\dots\le s+(s^2/2)/(1-s)\le s+s^2$.

\emph{Euler's constant} (first certificate only). With $H_n$ the $n$th harmonic number,
$\gamma=H_n-\ln(n+\tfrac12)-S_n$ with $0\le S_n\le2/(45n^2)$, and the script takes $n=2^{16}$
and sums $H_n$ in interval arithmetic. Indeed $c_k=H_k-\ln(k+\tfrac12)\to\gamma$, and with
$y=1/(2k)$, $c_{k-1}-c_k=\ln\frac{1+y}{1-y}-2y=2\sum_{j\ge1}y^{2j+1}/(2j+1)$, which lies
between $0$ and $(2y^3/3)/(1-y^2)\le(16/15)/(12k^3)$ for $k\ge2$; and
$\sum_{k>n}k^{-3}\le1/(2n^2)$. The result is
$\gamma\in[0.57721566490088,\ 0.57721566491124]$.

In tests at $80$ bits against $600$-bit reference values, the verified exponential and
logarithm showed no failure in $100{,}000$ points and $60{,}000$ wide intervals each. These
tests are evidence that the replacements are implemented as described; their correctness is
the argument above together with (B1).

\subsection{Certificate I: the integer constant}\label{sec:certificates-I}

\subsubsection*{The inequality}
The functional $\Psi_0(K,B,\theta,u)$ is defined in Definition~\ref{def:iid-psi0}. Let $K=5$,
$B=50$, $\theta=1.6712$, and let $u^\ast\in\R^{51}$ be the weight vector listed in
Proposition~\ref{prop:iid-certificate}; it is stored as five-decimal numbers in
\path{lemmaC/round3/firstmoment/w_K5_e10.json}. All inputs are exact rationals. The script
\path{lemmaC/round3/final/certI/certify_rigorous.py} certifies the inequality of
Proposition~\ref{prop:iid-certificate}:
\begin{equation}\label{eq:certificates-psi0bound}
\Psi_0(5,50,1.6712,u^\ast)\ \le\ -\frac{1162839537350571633950917}{2^{78}}\ \le\
-3.84751328322 .
\end{equation}
Consequently $-\Psi_0/\theta\ge2.30224586$, and
\begin{equation}\label{eq:certificates-exponents}
\theta\beta+\Psi_0\le
\begin{cases}
-0.0037532832,&\beta=2.30,\\
-0.0004108832,&\beta=2.302,\\
-0.0000766432,&\beta=2.3022,\\
-0.0000097952,&\beta=2.30224.
\end{cases}
\end{equation}
The second line is the criterion \eqref{eq:iid-criterion} of
Proposition~\ref{prop:iid-criterion} for $\beta=\betaint=2.302$, which gives
Theorem~\ref{thm:iid}. The certified threshold is $-\Psi_0/\theta\ge2.30224586$; the criterion
fails at $\beta=2.30225$.

For the market, Theorem~\ref{thm:integer} needs the exponent to stay negative after an
explicit error term $\mathrm{Err}$ is added, which accounts for the enlargement of the
typical event in the comparison (Section~\ref{sec:transfer},
Proposition~\ref{prop:transfer-interface}).
The script \path{constants_rigorous.py} in the same folder takes the eight-decimal
consequence $\Psi_0\le-3.84751328$ of \eqref{eq:certificates-psi0bound}, asserts it in exact
rational arithmetic against the exported endpoint, and certifies
\begin{equation}\label{eq:certificates-perturbed}
\theta\beta+\Psi_0+\mathrm{Err}\le
\begin{cases}
-0.0005929830\quad(\mathrm{Err}\le0.0031602970),&\beta=2.30,\\
-0.0001648848\quad(\mathrm{Err}\le0.0002459952),&\beta=2.302,
\end{cases}
\end{equation}
for the two parameter rows of the error budget (in the script's names: $c'=14$,
$\delta_1=10^{-5}$, $\delta_2=1/5000$, $\tau=1/200000$ for $\beta=2.30$, and $c'=17$,
$\delta_1=10^{-6}$, $\delta_2=1/50000$, $\tau=1/4000000$ for $\beta=2.302$).

\subsubsection*{The scheme}
Recall that $W=w_\infty+\sum_bD_bw_b$ with independent $D_b\sim\mathrm{Poisson}(\lambda_b)$,
$w_b=e^{u_b}$ and $w_\infty=e^{u_\infty}$. The Poisson expectations are carried exactly by the
Laplace transform
\[
\psi(t)=\E e^{-tW}=\exp\Bigl(-tw_\infty-\sum_b\lambda_b\bigl(1-e^{-tw_b}\bigr)\Bigr),
\]
so no truncation of the star size is needed. By Frullani's formula, with $\gamma$ Euler's
constant,
\begin{align*}
\Phi(u)=\E\ln W&=-\gamma+\int_0^1\bigl(1-\psi(t)\bigr)\frac{\dd t}{t}
-\int_1^\infty\psi(t)\frac{\dd t}{t},\\
\rho_b(u)&=\int_0^\infty w_be^{-tw_b}\psi(t)\dd t,\qquad
\eps_0(u)=\int_0^\infty w_\infty\psi(t)\dd t,
\end{align*}
and $\lambda_b$, $\kappa_b$, $J$ have closed forms. The functions $\psi$ and
$e^{-tw_b}\psi$ are Laplace transforms of positive variables, hence convex and decreasing.
The integrals are enclosed on a geometric mesh $t_0=2^{-27}$, $t_i=1.005\,t_{i-1}$ with
$5{,}109$ cells, reaching $t\approx868$. On each cell a convex function lies between its
tangent at the midpoint and its chord, which gives the midpoint and trapezoid rules as lower
and upper bounds for the $\dd t$-integrals and the corresponding closed-form bounds for the
$\dd t/t$-integrals. The head satisfies $0\le\int_0^{t_0}(1-\psi)\dd t/t\le t_0\,\E W$, and
for $t\ge t_1$ (the last node) $\psi(t)\le\psi(t_1)e^{-(t-t_1)w_\infty}$, which bounds the
tails; the mesh ends where these bounds are below $10^{-12}$. Finally,
$r\mapsto(1-r)^3+Cr^3$ and $e\mapsto-3e+Ce^3$ are convex, so their maximum over an enclosing
interval is at an endpoint, and the script takes the larger endpoint value.

\subsubsection*{Interval arithmetic and floating point}
Every quantity entering \eqref{eq:certificates-psi0bound}--\eqref{eq:certificates-perturbed}
is computed in interval arithmetic at $80$ bits, with (B1), (B2) and the verified
exponential, logarithm and Euler constant of Section~\ref{sec:certificates-defect}. The
weights are parsed as exact decimals. Powers are written as products. Floating point occurs in
three places, none of which affects validity: the mesh nodes are generated in floating point
and then used as exact binary rationals (the script asserts that they increase); the guess
$y_0$ in the verified logarithm; and the earlier search that produced the weights $u$, which
is irrelevant because the theorem holds for every $u$. Printed bounds are produced from the
exact binary endpoint by floor or ceiling at the stated number of decimals, and the
inequality is asserted. The exact endpoint is exported as a rational, with the parameters and
the SHA-256 hash of the weights file, to \path{psi0_export_w_K5_e10.json}; the second script
reads it and checks the hash.

\subsubsection*{Reproduction}
From \path{lemmaC/round3/final/certI/}:
{\footnotesize
\begin{verbatim}
python3 certify_rigorous.py ../../firstmoment/w_K5_e10.json 0.005   # 75 s
python3 constants_rigorous.py                                       # 1 s
python3 constants_rigorous.py --fresh                               # 80 s
python3 test_defect_and_replacements.py 100000                      # 4 min
\end{verbatim}}
The first command should print, among other lines,
{\footnotesize
\begin{verbatim}
Psi0 <= -3.84751328322   (8 decimals: Psi0 <= -3.84751328)
beta_cert = -Psi0/theta >= 2.30224586119   (8 decimals: >= 2.30224586)
  beta=2.302: theta*beta + Psi0 <= -0.0004108832  -> NEGATIVE (certified)
\end{verbatim}}
and the second, for $\beta=2.302$,
{\footnotesize
\begin{verbatim}
   theta*beta+Psi0 <= -0.0004108799 (m >= 0.0004108799);
   theta*beta+Psi0+Err <= -0.0001648848 (m - Err >= 0.0001648848)
   -> NEGATIVE (certified)
\end{verbatim}}
(one line in the output; broken here). With \texttt{--fresh} the second script recomputes
the endpoint in the same process instead of reading the export; both modes gave identical
output. The last command reproduces the counts of Section~\ref{sec:certificates-defect}.

\subsubsection*{Independent recomputations}
\begin{itemize}
\item A second certificate with other parameters ($K=5$, $B=20$, $\theta=1.6671$, weights
\path{w_K5_e25.json}) gives $\Psi_0\le-3.83621828482$ and $-\Psi_0/\theta\ge2.30113267639$.
\item A checking agent wrote its own interval enclosure and obtained $\beta\ge2.302028$, so
that the exponent at $\beta=2.30$ is at most $-0.00339$. This is weaker than
\eqref{eq:certificates-psi0bound} but independent. It still gives Theorem~\ref{thm:iid} with
$\betaint=2.302$, and in the market it covers the row $\beta=2.30$ of
\eqref{eq:certificates-perturbed} ($0.00339>0.0031603$) but not the row $\beta=2.302$.
\item A second checking agent read \path{certify_rigorous.py} line by line, ran it again,
and recomputed $\Psi_0=-3.847516394$ by two other methods (floating point).
\item Two checking agents recomputed the perturbed exponent at $\beta=2.30$ in exact rational
arithmetic and obtained at most $-0.0005929831$.
\item A reader from a different model family ran the original certificate (before the
replacements) and confirmed by direct interval comparison that $\Psi_0<-3.84751328$; it read
the logic of the hardened scripts without running them.
\item The original and the hardened scripts agree: the hardened bound is weaker by
$1.94\cdot10^{-12}$, all of it from the wider enclosure of Euler's constant, which enters
$\Psi_0$ with weight $3$.
\end{itemize}
Not done: a re-implementation in a different interval library.

\subsection{Certificate II: the cost of the tree matching}\label{sec:certificates-II}

\subsubsection*{The inequality}
The mathematics of this step is in Section~\ref{subsec:existence-rounding}, and we use its
notation: $\mathcal{Q}$ is the class of Definition~\ref{def:existence-Q}; for
$Q\in\mathcal{Q}$, $g_Q(t)=\E(t+X'+X'')_+^3/6$ with $X'$, $X''$ independent copies of $Q(E)$;
and $c(Q)=3\,\Gamma(Q)=3\bigl(\E\,\RDE(Q)(E)+2\,\E\,Q(E)\bigr)$ as in
\eqref{eq:existence-Gamma}, so that $c(Q^*)=9\,\E X$ at a fixed point. Rounded laws are
defined in Definition~\ref{def:existence-rounded}. The script
\path{lemmaC/round3/cert/cert.py} certifies the statement of
Proposition~\ref{prop:existence-certificate}: it produces rounded laws
$\tilde Q_0=0,\tilde Q_1,\dots,\tilde Q_{10}$ such that, for $0\le n\le9$ and
$g=g_{\tilde Q_n}$, the parameters $(\theta_i,\hat s_i)$ of $\tilde Q_{n+1}$ satisfy
\begin{equation}\label{eq:certificates-onesided}
\theta_i\le g(t_i),\qquad \hat s_i\ge g'(t_i)\qquad\text{at every mesh point } t_i,
\end{equation}
with $\theta_i$ strictly increasing and $\hat s_i$ nondecreasing, and it certifies
\begin{equation}\label{eq:certificates-cost}
c(\tilde Q_{10})\ \le\ 2.2519 .
\end{equation}
By Lemmas~\ref{lem:existence-rounding} and~\ref{lem:existence-chain}, the one-sided
inequalities \eqref{eq:certificates-onesided} give $\tilde Q_n\disp Q^*$ for every fixed
point $Q^*$, and then \eqref{eq:certificates-cost} gives $9\,\E X\le2.2519=\alphalp$, which
is Theorem~\ref{thm:existence}(ii). Table~\ref{tab:certificates-stages} lists the bounds
printed at all stages of the two recorded runs; only the last entry of run~B is used.

\begin{table}[ht]
\caption{Certified upper bounds on $c(\tilde Q_n)$, hence on the cost of every fixed
point, as printed (run~A: $h=2^{-10}$, $4000$ levels; run~B: $h=2^{-12}$, $12000$ levels),
and the exact iterates
$c_n=c(\RDE^n(\delta_0))$ in floating point for comparison. The floating-point limit is
$2.24954561$.}\label{tab:certificates-stages}
\begin{tabular}{@{}rccc@{}}
\toprule
$n$ & run A & run B & $c_n$ (floating point) \\
\midrule
0 & 4.867956 & 4.867955 & 4.86795 \\
1 & 3.417530 & 3.415996 & 3.41528213 \\
2 & 2.665068 & 2.661760 & 2.66021540 \\
3 & 2.404165 & 2.400021 & 2.39808802 \\
4 & 2.308352 & 2.303815 & 2.30169894 \\
5 & 2.274755 & 2.270053 & 2.26786005 \\
6 & 2.262954 & 2.258185 & 2.25596056 \\
7 & 2.258825 & 2.254029 & 2.25179199 \\
8 & 2.257381 & 2.252574 & 2.25033206 \\
9 & & 2.252065 & 2.24982093 \\
10 & & 2.251887 & 2.24964199 \\
\bottomrule
\end{tabular}
\end{table}

\subsubsection*{The scheme}
All bounds come from Lemma~\ref{lem:existence-enclosures}. Fix the lattice width $h=2^{-H}$
($H=10$ in run~A, $H=12$ in run~B). For the current rounded law, with cell masses $p_k$,
cell moments $D_k$ and chord weights $\nu_j$ as defined there:
\begin{itemize}
\item part (a) of the lemma, with termwise lower bounds
$p^{\mathrm{lo}}$, $D^{\mathrm{lo}}$, gives lower bounds $\theta_i\le g(t_i)$ and
$a_i\le g'(t_i)$ as sums over the lattice;
\item part (b), with termwise upper bounds $\nu^{\mathrm{hi}}$, gives upper bounds
$\hat s_i\ge g'(t_i)$;
\item part (c) gives an upper bound on $\E\,\RDE(\tilde Q)(E)$ from the $\theta_i$ and
$a_i$, and the exact formula for $\E\,\tilde Q(E)$ in terms of the cell integrals.
\end{itemize}
The script replaces $\hat s_i$ by the running maximum $\max_{k\le i}\hat s_k$, which keeps
\eqref{eq:certificates-onesided} and makes the slopes monotone; in both runs this changed no
entry. No tail estimate is needed: below $\theta_0$ the slope bound at $\theta_0$ is valid
for all smaller levels, and above $\theta_N$ the rounded slope is $0$. Both one-sided bounds
are of second order in $h$. The mesh points $t_i$ are multiples of $2^{-20}$, placed so that
$g(t_i)$ is roughly log-spaced on $[10^{-6},32]$. The printed number at stage $n$ is the
internal upper bound on $c(\tilde Q_n)$ rounded to nearest at six decimals
(Remark~\ref{rem:certificates-digits}).

\subsubsection*{Interval arithmetic and floating point}
\begin{itemize}
\item \emph{Interval arithmetic at $120$ bits} ((B1) and the verified exponential): everything
about a rounded law (breakpoints, the order comparisons between breakpoints and lattice
points, survival values, cell integrals, the mean) and the bound of
Lemma~\ref{lem:existence-enclosures}(c) with the final combination. Double-precision numbers enter
as exact rationals. An order comparison that the enclosures do not decide stops the program.
Conversion from an interval to a double uses the neighbouring double in the safe direction
and is asserted exactly.
\item \emph{Double precision with a proved error bound} ((B3)): the convolutions and the sums
of Lemma~\ref{lem:existence-enclosures}(a),(b). Every such quantity
is a sum of at most $2^{17}$ products of nonnegative numbers; the arguments $t_i+mh$ are exact
because both terms are multiples of $2^{-20}$ of size below $32$ (asserted); each term
involves at most three rounded multiplications. For such a sum, evaluated in any order and
with or without fused multiply-add,
$|\text{computed}-\text{exact}|\le((1+2^{-53})^{n+3}-1)\cdot\text{exact}+n\eta$ with
$n\le2^{17}$ terms and $\eta\le2^{-1022}$ for underflow. The relative error is below
$3\cdot10^{-11}$ and the absolute part below $3\cdot10^{-303}$. Every lower bound is
multiplied by $1-10^{-9}$ and reduced by $10^{-290}$ (and clamped at $0$); every upper bound
is multiplied by $1+10^{-9}$ and increased by $10^{-290}$. The margin also covers the
rounding of these adjustments and of the divisions by $2$ and $6$.
\item \emph{Plain floating point}: only the choice of the mesh points $t_i$, which has no
effect on validity, since Lemmas~\ref{lem:existence-rounding}
and~\ref{lem:existence-enclosures} hold for any mesh.
\end{itemize}

\begin{remark}[The quoted digits]\label{rem:certificates-digits}
The script prints its final bound with a round-to-nearest format, so the last printed digit
of $2.251887$ is not itself certified. Any internal bound that prints as $2.251887$ is below
$2.2518875$, and the bound quoted in this paper is $2.2519$. Three improvements are owed and
have not been made: asserting the internal bound against an exact rational, as in
Certificate~I; making the script refuse lattice parameters other than the two for which the
exactness argument above is written ($h=2^{-10}$ and $h=2^{-12}$); and saving the stage
parameters of run~B with the script's own output (regenerated copies are in
\path{lemmaC/round3/v10/}). \todo{make these three changes to cert.py and re-run}
\end{remark}

\subsubsection*{Reproduction}
From \path{lemmaC/round3/cert/}:
{\footnotesize
\begin{verbatim}
python3 cert.py 10 1500 2500 8 stages/A   # run A, under 1 min
python3 cert.py 12 4000 8000 10           # run B, about 4 min
python3 check_float.py stages/A 8         # 7 s
python3 test_small.py 4                   # 5 s
python3 test_rexp.py                      # 15 s
\end{verbatim}}
The arguments of \texttt{cert.py} are $H$, two mesh sizes, the number of stages and an
optional prefix for saving the stage parameters. Run~A should end with a line containing
{\footnotesize
\begin{verbatim}
stage 8: ... CERTIFIED c* <= 3*Gamma(Q~_8) <= 2.257381
\end{verbatim}}
and run~B with
{\footnotesize
\begin{verbatim}
stage 10: ... CERTIFIED c* <= 3*Gamma(Q~_10) <= 2.251887
\end{verbatim}}
The logs of the recorded runs are \path{run_A_h10.log} and \path{run_B_h12.log}. The mesh
choice is deterministic but may differ between builds of the linear-algebra library; any mesh
is valid, and the parameters of run~A are saved in \path{stages/}.

\subsubsection*{Independent recomputations}
\begin{itemize}
\item A checking agent wrote its own code, recomputed $g$ and $g'$ at all $4{,}000$ points of
all transitions of run~A (finding $\theta_i$ below $g$ and $\hat s_i$ above $g'$ everywhere),
obtained $c=2.25188597$ for the last law of run~B, and reproduced both runs.
\item A second checking agent computed the exact iterates $c_1,\dots,c_{10}$ with its own
code, in agreement with the values of Table~\ref{tab:certificates-stages} to all digits
shown, and reproduced both certificate runs digit for digit.
\item The script \path{check_float.py} is an independent code path in plain floating point
(cells of width $0.0007$ that are not dyadic, mass at cell midpoints, convolution by fast
Fourier transform, trapezoid rule). On the saved laws of run~A it gives
$c(\tilde Q_n)=4.867955$, $3.417527$, $2.665064$, $2.404162$, $2.308348$, $2.274751$,
$2.262951$, $2.258822$, $2.257377$ for $n=0,\dots,8$; the certified bounds exceed these by at
most $4\cdot10^{-6}$, and the certified parameters lie on the safe side of the floating-point
values at every mesh point.
\item The script \path{test_small.py} compares the bounds of
Lemma~\ref{lem:existence-enclosures}(a),(b) with tensor
Gauss--Legendre quadrature of the exact law on a $20$-piece example for $h=1/8$, $1/16$,
$1/32$: they are one-sided as required and of second order in $h$.
\item The certified values exceed the exact iterates $c_n$ at every stage, by $0.0007$ to
$0.0022$ in run~B, and the loss shrinks with the mesh, as expected of a rounding that only
loses dispersion. The final bound is $0.0023$ above the floating-point limit $2.24954561$ and
$0.050$ below $\betaint$.
\item A reader from a different model family audited the source of the script without
running it and found the enclosure scheme sound for both runs; the remark on the quoted
digits and the restriction on lattice parameters come from that audit.
\end{itemize}
An earlier version of the script used \texttt{iv.exp}; its logs are identical to those of the
present version.
